\section{Initial computational results}
\IfFileExists{generated/hubble_comparison.png}{
\begin{figure}[ht]
\centering
\includegraphics[width=0.82\linewidth]{generated/hubble_comparison.png}
\caption{Background expansion functions generated on a common redshift grid for standard
comparison models. Parameter values are illustrative configuration inputs, not fitted results.}
\label{fig:hubble-comparison}
\end{figure}}{}

\IfFileExists{generated/desi_dr2_bao_bbn_posteriors.png}{
\begin{figure}[p]
\centering
\includegraphics[width=\linewidth]{generated/desi_dr2_bao_bbn_posteriors.png}
\caption{DESI DR2 2025 BAO and BBN constraints in flat $\Lambda$CDM. The upper panels are
marginalized 68\% and 95\% posterior regions from the official released DESI MCMC chains:
BAO-only in the $(\Omega_m,h\,r_d)$ plane and BBN-calibrated BAO in $(\Omega_m,H_0)$. The
lower-left panel is the posterior-implied credible band for the model expansion history over the
DESI BAO redshift lever arm; it is not a continuous measurement of $H(z)$. The lower-right panel
is a CAMB background calculation showing $\Omega_{cb}$, $\Omega_\gamma$, total $\Omega_\nu$,
and $\Omega_\Lambda$, so the massive neutrino is not forced to scale as nonrelativistic matter at
early times. The analysis uses the pinned Results I/II likelihood; the later 2026 Lyman-$\alpha$
AP+BAO update is not combined without a release-native joint product.}
\label{fig:desi-dr2-bbn-posteriors}
\end{figure}}{}

\IfFileExists{generated/scpc_baseline_background.png}{
\begin{figure}[ht]
\centering
\includegraphics[width=0.9\linewidth]{generated/scpc_baseline_background.png}
\caption{Canonical \SCPC{} baseline trajectory and Friedmann-constraint diagnostic. This run
validates the pipeline and is not asserted to be cyclic.}
\label{fig:scpc-baseline}
\end{figure}}{}

\IfFileExists{generated/kinematic_diagnostics.png}{
\begin{figure}[ht]
\centering
\includegraphics[width=0.72\linewidth]{generated/kinematic_diagnostics.png}
\caption{Supplementary flat-$\Lambda$CDM kinematics. The equation-of-state curve is explicitly
$w_{\rm tot}=p_{\rm tot}/\rho_{\rm tot}$ rather than a dark-energy equation of state. The reference
lines $q=0$ and $w_{\rm tot}=-1/3$ identify the same acceleration transition in flat GR, which is
why this diagnostic is supplementary rather than a principal inference panel.}
\label{fig:kinematic-diagnostics}
\end{figure}}{}

\IfFileExists{generated/sound_horizon_backend_validation.png}{
\begin{figure}[ht]
\centering
\includegraphics[width=0.72\linewidth]{generated/sound_horizon_backend_validation.png}
\caption{Supplementary numerical validation of the Aubourg et al. sound-horizon approximation
against CAMB on BBN-calibrated posterior draws. The approximation is a cross-check only and is not
the inference backend.}
\label{fig:sound-horizon-validation}
\end{figure}}{}

\todo{Replace pipeline-validation figures with converged parameter scans, event classifications,
and stability results once a recurrent SCPC background has passed those gates.}
